% paper/sections/10_grid.tex
\section{界面適合非一様格子：追随格子幾何から高次演算子へ}
\label{sec:grid}

\noindent\textbf{本章の位置付け：}
前章までの圧力ジャンプ閉包と Balanced--Force 条件は，同じ界面面幾何を
圧力勾配，表面張力，補正速度が共有することを要求した．
非一様格子を導入する目的は，各時刻の界面帯と壁近傍へ格子点を配分し直し，
圧力勾配，表面張力，補正速度が共有する界面面幾何を現在格子上で再構築することである
（解析的格子伸長関数の体系は~\cite{Vinokur1983}）．
本章では，固定時刻の界面から定まる界面適合格子をまず定義し，
物理時間ステップごとの再生成では同じ構成を現在の界面へ適用する．
その各時刻の格子上で，座標変換，CCD/FCCD，Ridge--Eikonal が同じ幾何量を参照する流れを構成する
（離散対称性保存の観点は~\cite{VerstappenVeldman2003}）．
一様計算空間 $\xi$ では高次演算子をそのまま定義し，
物理空間 $x$ では連鎖律
（式~\eqref{eq:transform_1st_correct}, \eqref{eq:transform_2nd_correct}）
によって同じ離散幾何へ写す．
この分離により，格子生成の求積精度と微分演算子の精度を混同せずに議論できる．

\begin{tcolorbox}[warnbox, title={非一様格子の適用範囲}]
\label{box:nonuniform_scope_contract}
本章が本文で確定するのは，
\textbf{固定 $\varepsilon_\Gamma$（CLS/Heaviside 界面厚み）を保った非一様格子上で，
CCD/FCCD，Ridge--Eikonal，PPE と速度補正が同じ格子幾何を参照するための離散受理条件}
である．
本稿で検証する非一様経路は，固定非一様格子上の静止 BF 対照試験と，
表面エネルギー仮想仕事に基づく圧力ジャンプ/HFE/面上補正上の有限時間
固定幅・局所幅診断である（§~\ref{sec:nonuniform_grid_ns}）．
体積，鋭界面相体積，表面測度は一様計算座標ではなく
現在時刻の物理セル体積・物理面距離・物理界面交差位置で評価する．
毎ステップ界面追随格子は，保存的再写像，速度再投影，履歴リセット，
HFE 延長データと界面面幾何の再構築，および壁面境界条件での
PPE と速度補正の面演算子統一が同時に成立する場合にだけ標準化できる
\emph{閉包条件}として述べる．
保存形共通フラックス経路では，この保存的再写像は
$q$ だけでなくセル質量 $m$ と運動量 $\bm{p}_m$ を同じ写像で移す
$(q,m,\bm{p}_m)$ 再写像でなければならない．
格子だけを現在界面へ追随させ，速度を節点補間で載せ替える経路は，
見かけ上滑らかでも運動量と質量の同一性を証明しないため，
採用経路では受理しない．
時間固定格子は，界面が初期精細化帯内に留まる静止・小変形問題での
対照検証経路である．
一方，移動界面で上記閉包を欠く簡略再生成は，比較・診断経路に限る．
\end{tcolorbox}

\noindent\textbf{本章の構成：}
§~\ref{sec:grid_density} では界面帯・壁帯をまとめた格子モニタを定義し，
§~\ref{sec:grid_coord}--\ref{sec:grid_gen} では座標写像と計量係数を導く．
§~\ref{sec:grid_2d}--\ref{sec:grid_ale} ではテンソル積格子と追随格子の閉包条件を明確にし，
§~\ref{sec:eps_spatial_variable} では固定 $\varepsilon_\Gamma$ 標準経路から外れる空間可変幅の制約を述べる．
続く §~\ref{sec:ccd_metric}--\ref{sec:ccd_mixed} で CCD の非一様化を述べ，
§~\ref{sec:fccd_nonuniform} と §~\ref{sec:ridge_eikonal_nu_grid} で
面勾配と再初期化を同じ格子幾何へ接続する．

\subsection{格子密度関数}
\label{sec:grid_density}

界面近傍の解像度を上げる最小構成は，水準集合値に連動した密度関数
\begin{equation}
  \omega(\phi) = 1 + (\alpha - 1)\,\varepsilon_g\,\delta_g^*(\phi)
\end{equation}
である．$\alpha > 1$ は界面帯の格子密度倍率，
$\delta_g^*(\phi)$ は界面近傍に集中した正規化スムーズデルタ関数である．
$\delta_g^*$ は $[\text{長さ}]^{-1}$ の次元を持つため，
$\varepsilon_g\,\delta_g^*$ により無次元化する．

壁面を持つ問題では，界面の曲率・物性ジャンプ層とは独立に，
非すべり壁のせん断層，接触線閉包，投影法の半コントロールボリューム行が
壁法線方向の短い長さスケールを作る．
したがって本章の格子密度は，界面項と非周期壁項を同時に含む
方向別モニタとして定義する：
\begin{align}
  M_a(s) &= 1 + A_{\Gamma,a} I_{\Gamma,a}(s) + A_{W,a} I_{W,a}(s),
  \label{eq:grid_composite_monitor} \\
  \int_0^{x_{a,i}} M_a(s)\,\mathrm{d}s
  &= \frac{i}{N_a-1}\int_0^{L_a} M_a(s)\,\mathrm{d}s,
  \qquad i=0,\ldots,N_a-1 .
  \label{eq:grid_composite_equidistribution}
\end{align}
ここで $a$ は座標方向，$A_{\Gamma,a}=\alpha_{\Gamma,a}-1$，
$A_{W,a}=\alpha_{W,a}-1$ である．界面指示関数と壁指示関数は
\begin{align}
  I_{\Gamma,a}(s)
    &= \exp\!\left[-\frac{\bar{\phi}_a(s)^2}{\varepsilon_{\Gamma,a}^2}\right],
  &
  \bar{\phi}_a(s)
    &= \min_{x_b,\ b\ne a}|\phi(x)|,
  \\
  I_{W,a}(s)
    &= \max_{b\in W_a}\exp\!\left[-\frac{d_b(s)^2}{\varepsilon_{W,a,b}^2}\right],
  &
  d_{\mathrm{lower}}(s)&=s,\quad
  d_{\mathrm{upper}}(s)=L_a-s .
\end{align}
$W_a$ は方向 $a$ に垂直な物理壁の集合であり，周期境界では空集合とする．
周期境界は物理壁ではなく同一視境界なので，一方側だけに壁層を入れると
並進対称性を破る．したがって壁項は壁境界を持つ方向にだけ作用し，
一様格子または周期方向では $A_{W,a}=0$ として界面モニタのみへ退化する．

壁幅 $\varepsilon_{W,a}$ は安定化用の任意パラメータではなく，
物理壁層または明示的な数値解像度条件で決める：
\begin{equation}
  \varepsilon_{W,a}
  =
  \max\!\left(
    N_{W,\mathrm{cells}}h_{a,\mathrm{ref}},
    C_\nu\delta_W,
    C_{\mathrm{contact}}\varepsilon_{\Gamma,a}
  \right).
  \label{eq:grid_wall_width_policy}
\end{equation}
ここで $\delta_W$ は Stokes 層や境界層解析から得られる壁長さである．
それが与えられない場合は $N_{W,\mathrm{cells}}\ge4$ を
CCD による計量係数評価が解像すべき最小幅として明示する．

\paragraph{格子密度デルタ関数 \texorpdfstring{$\delta_g^*(\phi)$}{delta*g(phi)} の定義}

$\delta_g^*(\phi)$ は界面（$\phi=0$）近傍でのみ大きな値を持ち，
全域積分が正規化条件 $\int_{-\infty}^{\infty}\delta_g^*(\phi)\,d\phi = 1$ を満たす関数である．
格子生成に適した具体的な定義として，ガウス型スムーズデルタ関数を用いる：

\begin{equation}
  \delta_g^*(\phi) =
  \frac{1}{\varepsilon_g\sqrt{\pi}}\exp\!\left(-\frac{\phi^2}{\varepsilon_g^2}\right)
  \label{eq:grid_delta}
\end{equation}
ここで $\varepsilon_g$ は界面近傍の格子細分化幅（有効半幅）を制御するパラメータであり，中点値は $\varepsilon_g\,\delta_g^*(0) = 1/\sqrt{\pi} \approx 0.564$ となる（式~\eqref{eq:grid_delta}への代入より）．

\textbf{正規化：} ガウス積分 $\int_{-\infty}^\infty e^{-\phi^2/\varepsilon_g^2}\,d\phi = \varepsilon_g\sqrt{\pi}$ より $\int_{-\infty}^\infty \delta_g^*\,d\phi = 1$（$\checkmark$）．

\textbf{パラメータの選択指針：}
\begin{itemize}
  \item \textbf{界面幅を物理幅だけで指定する場合：}
        $\varepsilon_g = C_{g} \varepsilon_\Gamma$（$\varepsilon_\Gamma = C_\varepsilon h$）と置くと，
        格子点数 $N$ の一様計算格子幅 $\Delta\xi=L/(N-1)$ に対して
        $\varepsilon_g = C_g C_\varepsilon \Delta\xi$ となり，
        $\xi$ 空間でのモニタ幅は
        \[
          W_\xi = \frac{\varepsilon_g}{\Delta\xi} = C_{g} \cdot C_\varepsilon
        \]
        で固定される．
        例えば $C_{g} = 2$，$C_\varepsilon = 0.5$ では $W_\xi = 1$ セルである．
        計量係数 $J_x$ とその導関数 $\partial J_x/\partial\xi$ を CCD で
        6 次精度評価するには，格子密度関数が $\xi$ 空間で
        十分な幅に渡って変化している必要があり，
        $W_\xi \ge 4$ セルを下限の設計条件とする．
        $W_\xi = 1$ では $J_x(\xi)$ が 1 セル以内に急峻なピークを持ち，
        CCD の広域ステンシルで解像できず高次精度が失われる．
  \item \textbf{計算空間セル数定義（推奨）：}
        $\varepsilon_g$ を $\xi$ 空間のセル数で直接指定する：
        \begin{equation}
          \varepsilon_g = N_{g,\text{cells}} \cdot \Delta\xi,
          \qquad \Delta\xi = \frac{L}{N-1}
          \label{eq:eps_g_xi}
        \end{equation}
	        $W_\xi = \varepsilon_g/\Delta\xi = N_{g,\text{cells}}$ となり，$N$ や $\alpha$ によらず
	        $\xi$ 空間幅が $N_{g,\text{cells}}$ セルに固定される．
	        推奨値は $N_{g,\text{cells}} = 4$．
\end{itemize}

\noindent\textbf{格子生成の検証条件．}
非一様格子を用いる離散化は，少なくとも
$W_\xi\ge4$，
$\max_i H_{i+1}/H_i$（および逆比）の上限，
CCD で評価した $J_x,\partial_\xi J_x$ の格子収束を確認対象とする．
台形則生成のまま高次格子収束を主張してはならず，
その場合の主張は「計量係数の評価は高次だが，座標生成は
§~\ref{sec:grid_metric_order_note} の $\Ord{\Delta\xi^2}$ 律速を受ける」に限定する．

\noindent\textbf{格子密度関数の数値的特性：}
\begin{itemize}
  \item $\phi = 0$ では $\varepsilon_g\,\delta_g^*(0) = 1/\sqrt{\pi} \approx 0.564$，界面での格子密度は $\omega(0) = 1 + (\alpha-1)/\sqrt{\pi}$（例：$\alpha = 4$ のとき $\omega(0) \approx 2.69$）．
  \item $|\phi| > 3\varepsilon_g$ の領域では $\varepsilon_g\,\delta_g^*(\phi) < e^{-9}/\sqrt{\pi} \approx 7\times10^{-5} \ll 1$：ほぼ一様格子に戻る．
\end{itemize}

\subsection{座標変換と計量係数}
\label{sec:grid_coord}

物理座標 $x$ と計算座標 $\xi$ の連鎖律を適用すると，CCD が評価する $\partial/\partial\xi$ と物理空間 $\partial/\partial x$ を結ぶ計量係数 $J_x = d\xi/dx$ が現れる．両者の関係は以下の通りである．

\begin{align}
 \pd{f}{x} &= J_x\,\pd{f}{\xi}
 \label{eq:transform_1st_correct} \\[6pt]
 \pdd{f}{x} &= J_x^2\,\pdd{f}{\xi} + J_x\,\pd{J_x}{\xi}\,\pd{f}{\xi}
 \label{eq:transform_2nd_correct}
\end{align}
第2式は，$\partial f/\partial x = J_x\,\partial f/\partial\xi$ を $x$ で再微分し $\partial/\partial x = J_x\,\partial/\partial\xi$ を代入すると $J_x\,\partial/\partial\xi(J_x\,\partial f/\partial\xi)$ となり，積の微分則を展開して得られる．
ここで $J_x = d\xi/dx$ は計量係数であり，CCD 法を用いて高精度に評価される（評価構成は §~\ref{sec:grid_gen} 段階5；CCD 理論は §~\ref{sec:CCD}）．

\subsection{格子点座標の構成}
\label{sec:grid_gen}

計量係数 $J_x = d\xi/dx$ を得るには，まず物理座標配列 $\{x_i\}$ を
格子密度関数 $\omega(\phi)$ から生成する必要がある．
$\omega(\phi) > 1$ となる界面近傍では $\Delta\tilde{x} = \Delta\xi/\omega < \Delta\xi$ となり物理格子幅が細くなり，
$\omega \approx 1$ の遠方領域では均一格子に戻る；正規化（段階4）により総長は常に $L$ を保つ．
以下に，台形則による累積和として座標列を構成する
（概念図：一様格子→密度関数→界面適合格子の変換過程は図~\ref{fig:grid_transform} 参照）．

\begin{tcolorbox}[resultbox, title={1次元格子点座標の構成}]
\textbf{与える量：} 格子数 $N$，領域長 $L$，初期水準集合値 $\{\phi^{(0)}_i\}$（一様格子上），密度倍率 $\alpha$

\textbf{段階1：} 一様計算格子 $\xi_i = iL/(N-1)$（$i=0,1,\ldots,N-1$）を定義する．

\textbf{段階2：} 各格子点で $\omega_i = \omega(\phi^{(0)}_i) = 1 + (\alpha-1)\,\varepsilon_g\,\delta_g^*(\phi^{(0)}_i)$ を計算する．

\textbf{段階3：} 台形則で累積和を計算し，物理座標の「未正規化」値を求める：
\[
  \tilde{x}_0 = 0, \qquad
  \tilde{x}_{i+1} = \tilde{x}_i + \frac{\Delta\xi}{2}\left(\frac{1}{\omega_i} + \frac{1}{\omega_{i+1}}\right),
  \quad \Delta\xi = \frac{L}{N-1}
\]

\textbf{段階4：} $\tilde{x}_{N-1}$ で割って領域長 $L$ に正規化する：
\[
  x_i = L \cdot \frac{\tilde{x}_i}{\tilde{x}_{N-1}}
\]

\textbf{段階5：} 計量係数を CCD 微分で評価する：
\[
  J_x\big|_i = \frac{1}{(dx/d\xi)_i} \;(= d\xi/dx), \qquad
  \pd{J_x}{\xi}\bigg|_i = -\frac{(d^2x/d\xi^2)_i}{(dx/d\xi)_i^2}
\]
すなわち，座標列 $\{x_i\}$ に対して計算空間微分
$(dx/d\xi)_i$ と $(d^2x/d\xi^2)_i$ を同一の高次演算子で評価する．
素朴な中心差分による $J_x$ 評価は $\Ord{h^2}$ に留まり CCD の高次精度を損なうため採用しない（後掲 \ref{box:grid_jx_accuracy}）．

\textbf{得られる量：} 物理座標配列 $\{x_i\}$ と計量係数 $\{J_x|_i\}$，$\{\partial J_x/\partial\xi|_i\}$．
\end{tcolorbox}

\paragraph{\texorpdfstring{CCD $\Ord{h^6}$ 精度の平滑性条件}{CCD O(h^6) 精度の平滑性条件}}
CCD は評価対象配列が $C^8$ 級（8階微分有界）であるときに $\Ord{h^6}$ 精度を保証する
（打ち切り誤差は $f^{(8)}_i$ に依存，式~\eqref{eq:CCD_TE_II}）．
本稿の座標配列 $x(\xi)$ は $\omega(\phi)$（ガウス型 $\delta_g^*$ の滑らかな関数）の累積積分から得るため，
連続写像としては $C^\infty$ 級である．
一方，段階3で台形則を用いる場合，座標値そのものは連続写像に対して
$\Ord{\Delta\xi^2}$ の求積誤差を含む．
したがって「連続写像 $x(\xi)$ の計量係数を $\Ord{h^6}$ で近似する」と主張するには，
高次求積または解析的な累積写像が必要である．
本稿の標準構成で保証できるのは，生成済みの離散写像 $\{x_i\}$ に対する
CCD 微分評価が低次中心差分より高精度・低分散である，という範囲に限られる．

\noindent\textbf{精度律速の注意．}
\label{sec:grid_metric_order_note}
以上より，\textbf{CCD 微分評価自体は $\Ord{h^6}$ だが，
座標生成の台形則による $\Ord{\Delta\xi^2}$ 求積誤差が全体精度を律速する}．
高次格子収束を厳密に主張する場合は，段階3をシンプソン則（$\Ord{\Delta\xi^4}$）
または解析的 $x(\xi) = \int_0^\xi \omega^{-1}(\phi(\xi'))\,\mathrm{d}\xi'$ に置換する．
本稿の標準構成で台形則を用いる場合，精度主張は「CCD 微分演算子の内在精度」と
「格子生成の全体精度」を区別して記述する．
台形則と高次求積の差は，$J_x,\partial_\xi J_x$ および代表的な滑らかな試験関数の
格子収束で別々に確認する．

\begin{tcolorbox}[warnbox, title={計量係数 $J_x$ の精度と $\partial J_x/\partial\xi$ の評価}]
\label{box:grid_jx_accuracy}
素朴な中心差分 $J_x|_i \approx 2\Delta\xi/(x_{i+1}-x_{i-1})$ は $\Ord{h^2}$ 精度しか持たず，
CCD の $\Ord{h^6}$ 精度を著しく損なうため本稿では採用しない（段階5参照）．
その理由は，座標変換を用いた2階微分（式~\eqref{eq:transform_2nd_correct}）：
\[
  \pdd{f}{x} = J_x^2 \pdd{f}{\xi} + J_x \pd{J_x}{\xi} \pd{f}{\xi}
\]
の右辺第2項に $\partial J_x/\partial \xi$ が含まれるためである．
$J_x$ が $\Ord{h^2}$ 差分で与えられていると，
$\partial J_x/\partial\xi$ をさらに差分で評価した際に精度は $\Ord{h}$ 以下に劣化し，
CCD の $\Ord{h^6}$ 空間精度が $\Ord{h}$ または $\Ord{h^2}$ に低下する．
\end{tcolorbox}

\textbf{本稿の離散化選択：}
格子計量係数は座標列 $\{x_i\}$ から CCD 微分で評価し，
低次中心差分で得た $J_x$ と混在させない．
格子収束確認では，座標生成の求積次数と CCD 計量評価の次数を分けて評価し，
段階5の計量評価が採用する空間演算子と同じ幾何精度を持つことを確認する．

\begin{figure}[htbp]
\centering
\begin{tikzpicture}[>=stealth, font=\small]
% --- (a) 一様格子 ---
\begin{scope}[shift={(0,0)}]
  \node[above] at (2.5,2.3) {\textbf{(a) 一様格子}};
  \draw[gray!30] (0,0) grid[step=0.5] (5,2);
  \draw[thick] (0,0) rectangle (5,2);
  % 界面（円弧）
  \draw[blue, ultra thick] (2.5,1) circle [radius=0.7];
  \node[blue] at (2.5,0.1) {$\Gamma$};
\end{scope}

% --- (b) 密度関数 ---
\begin{scope}[shift={(6.5,0)}]
  \node[above] at (2,2.3) {\textbf{(b) 密度関数 $\omega(\phi)$}};
  \draw[->] (0,0) -- (4.2,0) node[right] {$\phi$};
  \draw[->] (0,0) -- (0,2.3) node[above] {$\omega(\phi)$};
  % 密度関数プロファイル
  \draw[red, ultra thick, smooth] plot[domain=0:4, samples=80]
    (\x, {0.4 + 1.6*exp(-(\x-2)^2/0.15)});
  \draw[dashed, gray] (0,0.4) -- (4,0.4) node[right, gray] {$1$};
  \draw[dashed, gray] (2,0) -- (2,2.0);
  \node[below] at (2,0) {$\Gamma$};
  \node[red, right] at (2.3,1.8) {$\omega(0)$};
\end{scope}

% --- (c) 非一様格子 ---
\begin{scope}[shift={(0,-3.5)}]
  \node[above] at (2.5,2.3) {\textbf{(c) 界面適合非一様格子}};
  % 非一様格子線（x方向：界面近傍で密）
  \foreach \x in {0, 0.45, 0.85, 1.2, 1.5, 1.7, 1.82, 1.92, 2.0, 2.08, 2.16, 2.24,
                   2.32, 2.42, 2.52, 2.64, 2.76, 2.88, 3.0, 3.1, 3.18, 3.3, 3.5, 3.8, 4.15, 4.55, 5.0}
    \draw[gray!50] (\x,0) -- (\x,2);
  % 非一様格子線（y方向：界面近傍で密）
  \foreach \y in {0, 0.25, 0.45, 0.58, 0.68, 0.76, 0.84, 0.92, 1.0, 1.08, 1.16, 1.24, 1.32, 1.42, 1.55, 1.75, 2.0}
    \draw[gray!50] (0,\y) -- (5,\y);
  \draw[thick] (0,0) rectangle (5,2);
  % 界面（円弧）
  \draw[blue, ultra thick] (2.5,1) circle [radius=0.7];
  \node[blue] at (2.5,0.1) {$\Gamma$};
  % 注釈矢印
  \draw[<-, thick, red!70!black] (2.5,1.8) -- (3.8,2.5) node[right] {格子集中};
\end{scope}
\end{tikzpicture}
\caption{界面適合非一様格子の生成過程}
\label{fig:grid_transform}
\PaperCaptionNote{(a) 初期の一様格子． (b) 界面 $\Gamma$ 近傍で密度関数 $\omega(\phi)$ がピーク値 $\omega(0)=1+(\alpha-1)/\sqrt{\pi}$ をとる． (c) 生成された非一様格子：界面付近で格子が自動的に集中する．}
\end{figure}
